% Resumo para o AERO2027 (World Summit on Aerospace and Aeronautical Engineering,
% Amsterdã, 24-26/06/2027). Prazo: 25/12/2026.
% Reproduz o modelo oficial (aero2027/modelo_oficial_AERO2027.docx):
%   página única A4, margens normais (2,54 cm), Times New Roman;
%   título 14 pt negrito centralizado, iniciais maiúsculas;
%   autores 12 pt (apresentador sublinhado; afiliação por sobrescrito; correspondente com *);
%   afiliações 11 pt itálico; texto 12 pt, entrelinha 1,5, justificado;
%   Keywords (máx. 6, 12 pt itálico); References; Biography (até 150 palavras).
% O modelo pede envio em Word: este .tex serve à redação no Overleaf; a versão final
% deve ser transposta para o .docx oficial sem alterar os estilos.
% Regras completas: deliverables/EVENTO_AERO2027.md e AGENTS.md.
% Trechos \pendente{...} (vermelho) ainda não têm fonte ou resultado. Não inventar números.
% Espaçamento: as entrelinhas simulam o "1.5 lines" do Word para Times (1,5 x 1,15 x corpo):
% 14 pt -> 24,15 pt; 12 pt -> 20,7 pt; 11 pt -> 18,98 pt. O texto do modelo pede A4 e 1,5,
% mas o arquivo .docx do modelo está em Carta e espaçamento simples (ver
% claude/outputs/D4_REVISAO_CLAUDE.md, R2); confirmar com os organizadores.
% Números: data/d4/contrastes_d4_5.csv e resumo_celulas_d4_5.csv (D4.1 v2.2, código 1262bb2).

\documentclass[12pt,a4paper]{article}
\usepackage[T1]{fontenc}
\usepackage[utf8]{inputenc}
\usepackage{mathptmx}                 % Times
\usepackage[margin=2.54cm]{geometry}  % margens "normais" do Word
\usepackage{xcolor}
\usepackage[normalem]{ulem}           % \uline para o apresentador
\usepackage{enumitem}

\newcommand{\pendente}[1]{\textcolor{red}{[#1]}}
\newcommand{\secao}[1]{\par\noindent{\fontsize{14}{24.15}\selectfont\bfseries #1}\par}
\pagestyle{empty}
\setlength{\parindent}{0pt}

\begin{document}

\begin{center}
{\fontsize{14}{24.15}\selectfont\bfseries
Coupling Terminal Station Queues And Stand Retention At Vertiports:\\
A Paired Simulation Of Admission Policy And Resource Capacity\par}
\medskip
{\fontsize{12}{20.7}\selectfont\bfseries
\uline{L. H. Tozi}\textsuperscript{1,*}, \pendente{coautores e professor}\textsuperscript{1}\par}
{\fontsize{11}{18.98}\selectfont\itshape
\textsuperscript{1}Instituto Tecnológico de Aeronáutica (ITA), \pendente{endereço}, Brazil.\par
\textsuperscript{*}\pendente{e-mail do autor correspondente}\par}
\end{center}

\secao{Abstract}
% texto justificado (padrão do LaTeX)
\fontsize{12}{20.7}\selectfont
Vertiport operators must decide whether scarce capacity should add a stand or a check-in counter,
and whether flights should be admitted in advance or first-come, first-served (FCFS). Airside
studies schedule pads and stands [1], terminal studies size passenger stations without the airside
[2], and coupled models report aggregate passenger delay rather than queues at each station [3].
We compare these decisions in a paired discrete-event simulation of a generic vertiport with one
TLOF and five stands, coupling charging and individual boarding to check-in, baggage and security
queues over thirty 8-h synthetic days. At 6 requests/h with 30-min charging, a fifth stand added
0.487 passenger departures/h under FCFS (95\% CI: 0.439--0.536) and 0.492 under advance admission
(0.440--0.544). A second counter cut mean check-in waiting by 141.5~s under FCFS (124.7--158.4)
and raised compliance with the IATA 180-s limit from 68\% to 99\%, but added 124.7~s
(107.9--141.5) of lounge waiting before boarding. Advance admission stayed within 0.033
departures/h of FCFS (0.001--0.066 fewer). The results compare designs under stated synthetic
inputs rather than predict a site.

{\fontsize{12}{20.7}\selectfont\itshape
\textbf{Keywords:} Vertiport; Urban Air Mobility; Discrete-Event Simulation; Passenger Processing;
Admission Policy; Level of Service\par}

\secao{References}
{\fontsize{12}{20.7}\selectfont
\begin{enumerate}[label={[\arabic*]},leftmargin=*,itemsep=0pt]
\item N. M. Guerreiro, G. E. Hagen, J. M. Maddalon, R. W. Butler, AIAA AVIATION 2020 Forum,
      doi:10.2514/6.2020-2903, (2020).
\item E. Feldhoff, AIAA AVIATION 2022 Forum, doi:10.2514/6.2022-3320, (2022).
\item L. Preis, M. Hornung, Electronics, \textbf{11}, 1071, (2022).
\end{enumerate}}

\secao{Biography}
{\fontsize{12}{20.7}\selectfont
L. H. Tozi is a master's student at Instituto Tecnológico de Aeronáutica (ITA), Brazil.
\pendente{completar com dados do autor; cabem cerca de 60 palavras no total sem passar de uma página com espaçamento 1,5}\par}

\end{document}
